\section{The arithmetic and geometric dictionary}
\label{sec:framework}

\subsection{Integral cohomology and rational mixed Tate motives}

Put $U_\ell=\operatorname{Spec}\mathbf Z[1/\ell]$.
For odd $\ell$, the exact comparison carried as K17-VK-PT is
\[
 K_{2e}(\mathbf Z)\otimes\mathbf Z_\ell
       \simeq H^2_{\acute et}(U_\ell,\mathbf Z_\ell(e+1)),
 \qquad e\ge2\text{ even}.
 \tag{FW.1}
\]
Its strict negative-twist dual is the Selmer group
\[
 \mathcal S_\ell(e)=\ker\left(
 H^1(U_\ell,\mathbf Q_\ell/\mathbf Z_\ell(-e))
 \longrightarrow
 H^1(\mathbf Q_\ell,\mathbf Q_\ell/\mathbf Z_\ell(-e))\right).
 \tag{FW.2}
\]
The relevant motivic weight is $e+1$: formally $2(e+1)-2=2e$.
Thus the completed new finite degrees concern twists $27$, $29$, $31$,
$33$, $35$ and $37$; all six have checked finite vanishing theorems.
The realizations in (FW.1), not merely their rationalizations,
retain the finite groups whose orders the Hecke certificates
control. This use of the modern integral edge comparison is
not an assertion that a factorial-normalized Chern character
is invertible at every prime.

Deligne--Goncharov, \emph{Groupes fondamentaux motiviques de
Tate mixte}, author manuscript \path{Tate.pdf}, \S1.6,
(1.6.11)--(1.6.12), printed p.~14, gives
\[
 \operatorname{Ext}^1_{\mathrm{MT}(F)}(\mathbf Q(0),\mathbf Q(n))
       =K_{2n-1}(F)\otimes\mathbf Q,\qquad
 \operatorname{Ext}^2_{\mathrm{MT}(F)}(\mathbf Q(0),\mathbf Q(n))=0.
\]
Their \S1.7 defines the unramified category over $\mathcal O_S$;
Proposition~1.9, printed p.~16, retains higher-weight
$\operatorname{Ext}^1$ and vanishing $\operatorname{Ext}^2$.
Proposition~2.3, printed p.~19, gives the free graded Lie algebra,
while \S2.4 distinguishes its canonical abelianization from
noncanonical free lifts. These are rational categories.
Consequently their second Ext vanishing cannot prove the
vanishing of the torsion in (FW.1). Nor is an integral abelian
category of mixed Tate motives with identical Ext formulas
being postulated here.

The coefficient sheaves in (FW.1) are continuous inverse limits
of the finite prime-invertible Tate sheaves. SGA~4, Expos\'e~IX,
3.2--3.3 is the finite Kummer/Hilbert~90 locator used in K17.
The convention for weights is Deligne, \emph{Weil~II},
1.2.1--1.2.2, p.~153, and 1.2.5(iv), p.~154:
$\mathbf Q_\ell(1)$ has weight $-2$ for geometric Frobenius.
The rank-two modular systems have pure weight $2e$ before
normalization; their twists in the trace comparison have weight
zero. The three trace terms are those of
$\mathbf Q_\ell(-e)\oplus\mathbf Q_\ell\oplus\mathbf Q_\ell(e)$
in the corresponding convention. Purity of rational systems
does not eliminate their integral congruences.

\paragraph{$\lambda$-rings and $L$-functions.}
SGA~6, Expos\'e~V, Definition~2.1, Example~2.2(b) and
Lemma~2.2.1 locate the $\lambda$-operations on $K_0$ defined
by exterior powers. The determinant and the first
$\gamma$-filtration piece record line bundles. The normal-line
Chern class in cote~149 is in this framework; its Kummer
identification is SGA~5, Expos\'e~VII, Proposition~3.4(i).
The higher groups in (FW.1) use the separately cited higher
$K$-theory comparison. No $K_0$ identity alone proves their
vanishing.

The CM formula in KT-GL2-CM describes the Euler factors of
rank-one Hecke characters induced from $E$. The checked endoscopic
construction decomposes a rank-three parameter
into a character and a rank-two base change. KT-TWOFIELD-661
uses the exact special values
$L(-2e,\chi_D)=-B_{2e+1,\chi_D}/(2e+1)$.
The rational class component requires the separate ordinary
Bernoulli number $B_{\ell-1-e}$; the two $L$-value questions
are not interchangeable. No new crystalline, prismatic, or
divided-power comparison is asserted.

\paragraph{The integral fixed locus and its conormal.}
For the source deformation algebra, the involution defines
the closed fixed locus by $Q=(s-\kappa_S(s))$.
Its conormal is $Q/Q^2$: this is the language of EGA~IV$_4$,
Definition~16.1.2, whose locator and notation are also identified
in the Stacks Project, Tag~\texttt{01R1}.  The affine formula
is Tag~\texttt{01R3}, and the base-change surjection, with an
isomorphism under flat base change, is Tag~\texttt{0473}.
These standard conormal statements do not assert flatness
of a fixed locus.  KT-UNSQUARED-NORMAL-JACOBIAN identifies
the required odd cotangent with its pullback to the base
augmentation.  KT-ODD-CONORMAL-SELMER-EQUIVALENCE then
identifies its Matlis dual with the actual strict Tate-sheaf
extension group in (FW.2), after extending coefficients to $O$.
Thus the formal geometry preserves the integral arithmetic
obstruction exactly.  The scalar plus normalized fourth
symmetric-power normal representation at an adjoint point
is the coefficient system displayed in KT-ADJOINT-ODD-NORMAL-MODULE;
it is not the Tate dual of that displayed coefficient system.
Generic automorphic point detection alone does not determine
the integral conormal or remove its vertical torsion.

\subsection{The Teichm\"uller tower and the letter to Faltings}

Grothendieck's \emph{Esquisse d'un programme}, \S2,
printed pp.~9--12 in \emph{Geometric Galois Actions~I},
describes the compatible Galois action on the profinite
fundamental groupoids of the curve-moduli tower, with gluing
and forgetting markings. Printed pp.~9--10 single out the
tripod and the first nontrivial moduli level; pp.~11--12
discuss generators, relations, and carefully chosen base-point
systems. Ihara's question here concerns the fixed field of
the kernel after selecting one pro-$p$ realization of the
tripod action. A full-profinite faithfulness statement is
not a proof of the claimed maximal pro-$p$ image, and a
rational unipotent image is another realization again.

The letter to Faltings, 27 June 1983, in the same volume,
pp.~49--58, makes the base-Galois augmentation explicit in
(1)--(3), pp.~50--52; p.~51 discusses successive hyperbolic
curve fibrations and Artin neighbourhoods. Formula~(7),
p.~54, is the point-to-section map. The correction on
pp.~56--57 treats cuspidal packets, and the postscript on
p.~58 corrects the open-image/dominance conditions and gives
a proposed cyclotomic fixed-subgroup criterion. These
locators were read in the published transcription.
The finite-field counterexample below concerns an extension
beyond the characteristic-zero section conjecture formulated
there. The higher-dimensional counterexamples concern a
specific maximal geometric prime-to-$p$ quotient; they do
not refute the letter's full profinite formulation.

\subsection{Cote and run crosswalk}

The following table specifies the object and implication for
each family of findings. These are applications or limitations
of the named frameworks, not claims that Grothendieck's texts
already contain the later numerical theorems.

\small
\begin{longtable}{>{\raggedright\arraybackslash}p{.24\textwidth}
>{\raggedright\arraybackslash}p{.34\textwidth}
>{\raggedright\arraybackslash}p{.34\textwidth}}
\toprule
Finding family & Geometric or arithmetic object & Exact programme locator and change\\
\midrule\endhead
KT-ONECHAR-CLOSED; KT-TRACE-MULTIPLICITY; KT-QUARTIC-BOUND &
Strict extension classes of negative Tate sheaves, dual to
integral $H^2(U_\ell,\mathbf Z_\ell(e+1))$. &
Alp\"oge step~17; K17-S17-UNIT and VK-PT. Bridges Q06,
GRH\_NIGHT AG-Q18: a bound on every $K_{4k}$ is now explicit,
while the all-weight zero statement remains open.\\
KT-U3-WEIGHT, DIM, SECOND, SPECTRUM; KT-U3-ENDO-EXACT; KT-U3-SPECTRAL-SPLIT; KT-GL2-CM &
Algebraic automorphic coefficient spaces and Frobenius Euler
factors, with cohomological weight shift retained. &
K17-U3-HECKE, U3-GENERAL-COVERAGE, GL2-GENERAL; prompt
items1(a)--(b),2. These give uniform finite models, not a
universal theorem on their residual primes.\\
KT-K52-K56-K60-VANISH; KT-K64-VANISH; KT-K68-VANISH; KT-K72-VANISH;
KT-COMPONENT-EXTENSION-18 &
Twists27,29,31,33,35,37 of arithmetic $H^2$, the cyclotomic even components
$[-26],[-28],[-30]$, $[-32],[-34],[-36]$, and the corresponding odd ramified modules. &
K17 Q2-HIGHER4; K17-COMPONENT-CONTROL; LM6 A.44--A.45;
bridges Q06. The character interval grows to $[-16,18]$.
No complete remaining-prime kernel assertion follows.\\
KT-C7-EXPLICIT; KT-SOULE-STAB &
Torsion cohomology of arithmetic groups and integral higher
$K$-theory of rings of integers. &
K17 VK-EXP7; Soul\'e2003 \S4, Theorem2 and Proposition4;
Lemma9 pp777--778 requires the stated general-field correction.\\
KT-TWOFIELD-661; KT-SCHUR-PIVOT-SUPPORT &
Actual Eisenstein congruences, reflected rational characters,
and units in finite coefficient rings. &
K17-REFLECTOR (RB.1); prompt1(b)--(c). They identify
limitations of proposed uniform shortcuts; neither gives
nonzero $K_{4k}$ or a Vandiver counterexample.\\
KT-U3-SYM2-INJECT; KT-ORDINARY-SYM2-ADM3;
KT-ORDINARY-SYM2-ALL-LEVELS;
KT-STABLE-PROPAGATION; KT-ADJOINT-ACTUAL-HOST-POINT &
Weight-zero rank-three lisse sheaves with ordinary flags,
residual Tate constituents, cuspidal automorphic lifts, and
actual points of the normalized deformation ring. &
K17-S17-UNIT, actual Adm3U and CoversPt sources; the integral
Tate/Selmer dictionary (FW.1) is retained. Their congruences
can coexist with zero integral motivic cohomology. The strict
Selmer tangent line has the opposite sign from the adjoint
root; the new conormal comparison below retains that sign.\\
KT-ODD-CONORMAL-GATE; KT-ODD-CONORMAL-SELMER-EQUIVALENCE;
KT-FINITE-GRADED-NORMAL-MODULE &
An integral fixed-locus conormal, its weight and involution
summands, and the strict negative Tate extension group. &
EGA IV$_4$,16.1.2; SGA4 IX,3.2--3.3 for the finite Tate-sheaf
language; actual source ModelLower and ShaInjGen.
The conormal is exactly the Selmer dual. A unit minor remains
equivalent to the arithmetic vanishing problem; finite
generation does not supply flatness or a unit determinant.\\
KT-STABLE-HECKE-SELFDUAL-LOCALIZATION;
KT-TRACE-INVISIBLE-SELMER-TANGENT;
KT-RESIDUAL-TRACE-DISCRIMINANT &
Integral Hecke orders, Frobenius characteristic polynomials,
and extension data lost by the trace map at the reducible point. &
Alp\"oge step17, ModelLower and ShaInjGen; the conormal
framework EGA IV$_4$,16.1.2 and (FW.2).
An evaluated complementary factor gives an all-split-prime
Hecke identity, while the actual Selmer tangent is invisible
to traces and the split-residual matrix trace order has
nonunit discriminant. No unit normal comparison follows.\\
KT-FIELD-CHARACTER-REDUCTION; KT-EXPLICIT-TWO-UNIT-FAMILY;
KT-BERNOULLI-PREFIX-FORMULA; KT-THIRD-UNIT-CHARACTER-OBSTRUCTION;
KT-AUXILIARY-UNIT-SELMER-BRIDGE; KT-RATIONAL-POSITIVE-ROOT-UNIT &
Special values of odd Dirichlet characters with finite
Teichm\"uller twists and a fixed negative discriminant. &
Prompt1(b)--(c), K17-REFLECTOR and the integral dictionary
(FW.1)--(FW.2). The exact same-discriminant three-unit
condition is not rational mixed-Tate Ext vanishing.
For the explicit weight progression (EU.2), two fields
supply its first two units and opposite local behavior
at the residue prime. The faithful cyclic-character projector
annihilates all units imported from multiquadratic fields; full
compositum $p$-rationality is stronger than the third-unit condition.
These are character decompositions of the finite Tate/Dirichlet
coefficient sheaves in (FW.1)--(FW.2), not vanishing of the
original rational strict extension group.\\
KT-GMA-REDUCIBILITY-DOES-NOT-BOUND-EXTENSIONS;
KT-RESIDUAL-CYCLICITY-SCOPE;
KT-POLARIZED-EXTENSION-GRAPH-OBSTRUCTION &
Integral extension modules of a residual sum of three
rank-one sheaves; involution signs and mixed compositions
of extensions. &
The conormal framework EGA IV$_4$,16.1.2 and the exact
strict coefficient sheaf in (FW.2). Trace and reducibility
ideals alone do not bound the needed extension module.
The abstract graph retains irreducibility after five full
cohomology vanishings; it is not an arithmetic realization
or a counterexample to motivic vanishing.\\
KT-ORDINARY-RANK2-UNIT-EXCLUSION;
KT-AUXILIARY-UNIT-NO-TYPEII &
An ordinary rank-two sheaf, its opposite integral lattice,
and the resulting strict extension of a quadratic Tate sheaf. &
Alp\"oge step17, actual Type-II points; source
\texttt{TanDet.fin\_quot} and \texttt{typeII\_facts}.
A quadratic Bernoulli unit excludes this source branch
without a weight bound. The remaining rank-three extension
problem is still an integral obstruction in (FW.2).\\
KT-QUARTIC-QUARTER &
A refined finite-order bound for integral motivic cohomology
of Spec Z. &
K17 VK-PT and VK-EXP7; Deligne--Goncharov1.6(1.6.11)--(1.6.12)
records the rational comparison, which forgets this torsion.\\
KT-FF-SEC &
A point-free proper curve over $\mathbf F_2$ with a split
arithmetic fundamental-group extension. &
Cote149 scans35--36, AN\_GRH A05. Point separation does not
imply section surjectivity. SGA1 IX.6.1 supplies the
fundamental-group exact sequence.\\
KT-POLY-FAIL; KT-ARTIN-FAIL; KT-PRODUCT-SEP &
Wild isotrivial twisting lost on maximal geometric
prime-to-two completion; separation retained by maps to curves. &
Extension of cote149 A1 and AN\_GRH A05; Faltings letter p51.
SGA1 X.1.7, XIII.4.6 are the proper and prime-to-characteristic
Kunneth inputs in the respective proofs.\\
KT-F149-REVIEW &
Geometric tangent vectors inside a completed Kummer torsor,
with the generic Galois marking. &
Cote149 scans83--84, generic diagram; scan85, geometric
subtorsor. K17 F149-GENERIC, MARKED-RECOGNITION and
UNMARKED-CURVE. EGA IV$_2$,5.10.5 supplies normal extension;
the required anabelian reconstruction is separately cited.\\
KT-F149-COMP &
The normalizer of cusp inertia, the nonzero-vector line bundle,
and its integral first Chern class. &
Cote149 scans79--82; GRH\_NIGHT AG-C149-NORMAL, FULLLIFT,
CUSP and K17 F149-RELATIVE. Normal/conormal signs are kept
consistent. SGA5 VII.3.4(i) locates the Chern/Kummer language.\\
KT-R2-POINT and its remaining incidence question &
Odd cyclic etale covers of an ordinary genus-two curve,
primitive two-rank, and the Raynaud theta divisor. &
AN\_GRH A33 and A45--A52; original failed constructions
A36--A44 are retained. Over $\overline{\mathbf F}_2$ pointing
is automatic, but the RM incidence problem remains open.\\
KT-RM-FROB-AVOID; KT-RM-FROB-NONDETERMINATION &
The finite etale torsion local systems of the universal
Hilbert--Blumenthal Jacobian and the actual relative
Raynaud theta support. &
AN\_GRH A19 (FS.3)--(FS.10), A37 Steps1--5;
cote149 scans35--36 and its route-two substitute A33.
Arbitrary finite symplectic Frobenius data coexist with
avoidance at all exponents of any fixed odd prime support.
The all-odd-order cover question remains open.\\
Ihara gate &
Actual arithmetic pro-$p$ action and its plus Iwasawa modules,
including the cubic companion errors. &
Esquisse2; LM6 A.44--A.45; K17 Q3-IRREG and CUBIC-ODD;
GRH\_NIGHT AG-Q29. Rational Brown/Deligne--Goncharov
freeness does not supply arithmetic Triple-NV.\\
\bottomrule
\end{longtable}
\normalsize

The FGA representability programme provides the language of
families, Picard functors, and moduli. No FGA theorem is
invoked as a missing anabelian reconstruction or integral
$K$-theory comparison in this report, so no spurious FGA
statement number is assigned to one of these new assertions.
